\documentclass[german,10pt]{book}      \input{preambel}         \begin{document}  \setcounter{chapter}{1}


\chapter{Observable und Zust\"ande}
% Kap x
\label{chap_Observable}

\info{Thomas Filk}{5.10.2025}%

In jeder physikalischen Theorie gibt es zwei Konzepte, auf denen die Theorie beruht:
(1) das Konzept der Observablen, und (2) das Konzept des Zustands. In der Quantentheorie
ergeben sich die teilweise erheblichen Unterschiede zwischen den verschiedenen Interpretationen oftmals
aus den unterschiedlichen Vorstellungen \"uber diese beiden Konzepte, insbesondere dem des
Zustands. In diesem Kurztext soll zumindest ein grober \"Uberblick \"uber die verschiedenen
Vorstellungen zum Zustandsbegriff in der Quantentheorie gegeben werden.

\section{Observable}

\textbf{Definition:} \textit{Observable sind physikalische Gr\"o\ss en, 
die man an einem physikalischen System messen kann. Gelegentlich versteht man unter
einer Observablen auch die mathematische Repr\"asentation dieser Gr\"o\ss en.}
\vspace{0.3cm}

\noindent
John von Neumann
definiert ein physikalisches System durch die Menge der Observablen, die an diesem
System beobachtet werden k\"onnen \cite{Neumann}. Welche Observablen das sind,
ist eine Erfahrungstatsache. Erst durch das Ergebnis des Stern-Gerlach-Experiments 1922
erkannte man, dass es neben Ort und Impuls noch einen weiteren Freiheitsgrad gibt, den man an einem
Elektron messen kann, n\"amlich seine Spin\-orien\-tierung entlang einer vorgegebenen Achse. 
Man kann daher nie behaupten, s\"amtliche an einem System messbaren Observablen zu kennen, bzw.\ es 
kann nie ausgeschlossen werden, dass in Zukunft noch weitere Observablen bekannt werden.
Ein allgemeiner physikalischer Formalismus muss f\"ur eine Theorie zun\"achst
kl\"aren, welche Observablen es gibt bzw.\ bekannt sind und wie diese Observablen mathematisch repr\"asentiert
werden sollen. Observablen bilden jedoch nicht einfach nur eine strukturlose Menge, sondern es gibt auch 
Beziehungen zwischen den Observablen, die bekannt sein sollten und sich in der mathematischen Darstellung
widerspiegeln m\"ussen.

Die physikalische Realisierung einer Observablen erfolgt durch die Angabe des Mess\-proto\-kolls, wie
eine Messung dieser Observablen durchzuf\"uhren ist. Es ist eine Abfolge von Schritten, die angeben, in
welcher Form ein System mit einer Messapparatur gekoppelt wird und wie die 
Ver\"anderungen an der Messapparatur (z.B.\ die Zeigerstellung) abgelesen werden. Aus diesen
Ver\"anderungen muss sich eine Zahl bestimmen lassen, die man als das Ergebnis der Messung bezeichnet
und die eine Aussage \"uber den Zustand eines Systems -- wie es pr\"apariert wurde -- macht.
In diesem Kapitel bezeichne ich die physikalische Observable
als Messvorschrift in Anlehnung an von Neumann \cite{Neumann} mit kaligraphischen Buchstaben, z.B.\
${\cal R}$, die zugeh\"orige mathematische Darstellung dieser Observablen kennzeichne ich 
durch $R$. Sehr oft wird in der Literatur leider nicht zwischen diesen beiden Objekten unterschieden und es wird
erwartet, dass man im Einzelfall entscheiden kann, ob die physikalische Messvorschrift oder
die mathematische Darstellung gemeint ist.

Kennt man eine Observable ${\cal R}$, also die Vorschrift, wie man diese Observable an einem System 
messen kann, so kennt man auch die Observable $f({\cal R})$ f\"ur eine Funktion $f$: Man misst die
Observable ${\cal R}$ an einem physikalischen System und bildet von dem erhaltenen Messwert $r$ die
Funktion $f(r)$. Insbesondere ist die Observable $\alpha {\cal R}$ (f\"ur eine beliebige reelle Zahl $\alpha$)
durch die Messvorschrift von ${\cal R}$ gegeben, wobei jeder gemessene (an einer Zeigerstellung
abgelesene Messwert) mit $\alpha$ multipliziert wird. Lassen sich zwei Observable ${\cal R}$ und ${\cal S}$
gleichzeitig an einem System messen (das bedeutet, es gibt ein Protokoll, bei dem man gleichzeitig
einen Messwert sowohl f\"ur ${\cal R}$ als auch f\"ur ${\cal S}$ erh\"alt -- diese Werte seien $r$ und $s$), 
dann ist 
die Observable $f({\cal R},{\cal S})$ definiert als die  Messvorschrift, bei der von den Messwerten die Funktion
$f(r,s)$ gebildet wird. 

Bei Observablen ${\cal R}$ und ${\cal S}$, die sich nicht gleichzeitig an einem System messen lassen -- solche
Observablen bezeichnet man als nicht kompatibel --, sind
allgemeine Funktionen $f({\cal R},{\cal S})$ zun\"achst nicht definiert. Insbesondere sind auch Ausdr\"ucke der Art
${\cal R}+{\cal S}$ oder ${\cal R}\cdot {\cal S}$ als Messvorschriften nicht definiert.   

Die mathematische Darstellung (Repr\"asentation) einer Observablen h\"angt von der Theorie
ab, die wir verwenden. In der klassischen Mechanik handelt es sich bei Observablen um 
Funktionen von Ort und Impuls (bzw.\ Geschwindigkeit), d.h.\ Funktionen auf dem
Phasenraum. In der Quantenmechanik
werden Observable durch sogenannte selbstadjungierte bzw.\ hermitesche Operatoren auf einem
Hilbert-Raum -- einem Vektorraum mit einem Skalarprodukt -- dargestellt. 

An dieser Stelle sollte man betonen,
dass die mathematische Darstellung einer Observablen nicht den Messprozess repr\"asentiert, also
den dynamischen Vorgang der Messung, sondern eher die Informationen, die man durch solche
Messungen erlangen kann: die Menge der m\"oglichen Messwerte sowie die Zust\"ande, die mit
diesen Messwerten verbunden sind. In diesem Sinne (und in Anlehnung an einen Ausdruck von
Schr\"odinger \cite{Schroedinger} zum Begriff der Wellenfunktion) 
kann man von einer Observablen als einem \glqq Katalog von
m\"oglichen Ergebnissen\grqq\ sprechen. 

\section{Zustand}

\subsection{Zustand als Erwartungswertfunktional}

Auch bei Zust\"anden sollte man zwischen der physikalischen Realisierung und der
mathematischen Darstellung unterscheiden. Sehr oft, insbesondere in mathematischen Texten
(z.B.\ dem Buch \glqq Mathematische Grundlagen der Quantenmechanik\grqq\ von John von
Neumann \cite{Neumann}) definiert man einen Zustand als
ein \glqq Erwartungswertfunktional auf der Menge der Observablen\grqq. Diese zun\"achst
mathematische Definition kann man durchaus w\"ortlich interpretieren: Wenn wir ein System
durch einen Zustand beschreiben, erlaubt es dieser Zustand, jeder Observablen eine Zahl -- ihren
Erwartungswert in diesem Zustand -- zuzuordnen. Ein Zustand beschreibt also unsere Erwartungen
in Bezug auf die Ergebnisse, die bei einer Messung an einem System auftreten. 
Schr\"odinger bezeichnet einen Zustand als einen \glqq Katalog der Erwartung\grqq\ \cite{Schroedinger}. 

Dieser Katalog der Erwartung beruht auf unserem Wissen 
\"uber die Vergangenheit eines Systems bzw.\ eines Ensembles von Systemen. 
Im Wesentlichen bezieht sich dieses Wissen auf die Art, wie diese Systeme pr\"apariert wurden. 
Unsere Erwartungen in Bezug auf m\"ogliche Messergebnisse k\"onnen nur auf diesem Wissen
beruhen, und wenn uns bewusst ist, dass unser Wissen unvollst\"andig ist, verwenden
wir sogenannte gemischte Zust\"ande (in der klassischen Mechanik sind das Wahrscheinlichkeitsverteilungen
\"uber einem Zustandsraum, beispielsweise dem Phasenraum; in der Quantentheorie sind das
sogenannte Dichtematrizen). Unser Wissen \"uber die Vergangenheit
eines Sys\-tems muss nicht beliebig weit zur\"uckreichen: Es reicht das Wissen \"uber solche Aspekte, die
f\"ur die Vorhersagen zuk\"unftiger Messungen von Bedeutung sind. Das sind meist die letzten
Pr\"aparationen, die an einem System vorgenommen wurden.\footnote{Die Physik ist hier in einer
gl\"ucklichen Lage. M\"ochte man in der Psychologie den mentalen Zustand einer Person beschreiben,
k\"onnen beliebig weit zur\"uckliegende Ereignisse oder Erfahrungen f\"ur die Vorhersage des
zuk\"unftigen Verhaltens wesentlich sein.} Es gibt aber auch physikalische Systeme mit 
einem \glqq Ged\"achtnis\grqq, beispielsweise Spinglassysteme oder neuronale Netzwerke. 

Wir k\"onnen einen Zustand somit definieren als eine Kodierung unserer Kenntnis \"uber
die Vergangenheit eines physikalischen Systems, soweit diese Kenntnis es uns erlaubt,
zumindest statistische Vorhersagen \"uber die Ergebnisse zuk\"unftiger Messungen (von Observablen)
zu treffen. 

Die physikalische Realisierung eines Zustands erfolgt durch ein im Prinzip beliebig gro\ss es
Ensemble gleichartig pr\"aparierter Systeme. Gew\"ohnlich f\"uhren wir ein Experiment nicht nur
einmal durch, sondern nach M\"oglichkeit unter denselben Ausgangsbedingungen gen\"ugend oft,
sodass wir statistische Auswerteverfahren verwenden k\"onnen. Au\ss erdem k\"onnen wir auf diese
Weise an demselben Zustand verschiedene Observablen messen, die sich an demselben System
nicht gleichzeitig bestimmen lassen (beispielsweise Ort und Impuls bei quantenmechanischen
Systemen oder allgemeiner zwei nicht kompatible Observablen ${\cal R}$ und ${\cal S}$). 
Man unterteilt dazu das Ensemble gleichartig pr\"aparierter Systeme in zwei
ebenfalls gro\ss e Teilensembles und misst an dem einen Teilensemble die Observable
${\cal R}$ und an dem anderen Teilensemble die Observable ${\cal S}$. Auf diese Weise erh\"alt man f\"ur
einen Zustand -- das Ensemble gleichartig pr\"aparierter Systeme -- die Erwartungswerte von
nicht kompatiblen Observablen in diesem Zustand.

Gelegentlich wollen wir auch einem Einzelsystem einen Zustand zuschreiben. An diesem
Punkt unterscheiden sich verschiedene Interpretationen, insbesondere die
sogenannten Ensemble-Interpretationen von der Kopenhagener Deutung. W\"ahrend
Ensemble-Interpretationen den Begriff der Wahrscheinlichkeit vermeiden und nur von
relativen H\"aufigkeiten sprechen -- die Born'sche Regel besagt dann, dass das
Absolutquadrat eines Skalarprodukts gleich der relativen H\"aufigkeit ist, mit der bei einer
Messung einer Observablen an einem Ensemble gleichartig pr\"aparierter Zust\"ande ein 
bestimmtes Resultat erzielt wird --
verwendet die Kopenhagener Interpretation den Zustandsbegriff auch f\"ur Einzelsys\-teme
und spricht statt von relativer H\"aufigkeit von Wahrscheinlichkeit. Es gibt aber auch
Situationen, bei denen der Unterschied relevant wird: Beispielsweise gilt das sogenannte
No-Cloning Theorem, wonach ein unbekannter Zustand nicht dupliziert werden kann, 
nicht f\"ur Ensemblezust\"ande. An einem Ensemble kann der Zustand bestimmt und dann
beliebig oft kopiert werden. Diese Aussage gilt jedoch nicht f\"ur den Zustand, den wir einem
einzelnen System zuschreiben. 

Man bezeichnet einen Zustand --
realisiert als Ensemble -- als rein, wenn es keine Unterteilung dieses Ensembles in (f\"ur statistische Auswertungen
ebenfalls gen\"ugend gro\ss e) Subensembles gibt, die f\"ur eines dieser Subensembles andere
Erwartungswerte liefert als das gesamte Ensemble. Diese Vorstellung von reinen Zust\"anden entspricht
auch unserer \"ublichen Vorstellung von reinen
Zust\"anden als maximale bzw.\ nicht mehr erweiterbare Kenntnis \"uber ein System. Entspricht
ein Ensemble von Systemen einem Gemisch, so l\"asst es sich in Subensembles aufteilen, die
weniger gemischten Zust\"anden (oder sogar reinen Zust\"anden) entsprechen. Reine Zust\"ande
lassen sich nicht mehr in noch reinere Zust\"ande aufteilen. 
 
Mathematisch wird ein Zustand repr\"asentiert durch die Angabe einer Abbildung, die
jeder Observablen eine Zahl -- ihren Erwartungswert in diesem Zustand -- zuordnet. 
Ein Zustand ist somit eine mathematische Kodierung unserer Kenntnis \"uber die Art, wie ein
System pr\"apariert wurde, sodass wir dieses Wissen f\"ur die Vorhersage zuk\"unftiger
Messungen -- die Vorhersage eines Erwartungswerts f\"ur beliebige Observable -- an diesem 
System nutzen k\"onnen. 

Elementare Eigenschaften dieser Abbildung sind: 
\begin{enumerate}
\item 
Die Identit\"atsobservable, die immer nur
den Messwert 1 als Ergebnis liefert, soll den Erwartungswert 1 haben, 
\item 
eine positive Observable,
die immer nur positive Ergebnisse als Messwerte liefert, soll auch einen positiven Erwartungswert
haben, und 
\item
das $\lambda$-fache einer Observablen ${\cal R}$, die immer das $\lambda$-fache eines
Messwerts von einer Messung von ${\cal R}$ liefert, soll auch das $\lambda$-fache des Erwartungswerts
haben. 
\end{enumerate}
Weitere Eigenschaften h\"angen davon ab, welche Strukturen auf der Menge der
Observablen definiert sind. 

Oft unterscheidet man (insbesondere in der Quantentheorie) den Zustand von dem Erwartungswertfunktional, 
das dieser Zustand definiert. Der Zustand wird in der Quantentheorie z.B.\ durch einen normierten
Vektor $|\psi\rangle$ in einem Hilbert-Raum dargestellt, wohingegen die Abbildung 
${\cal R} \rightarrow {\rm Erw}({\cal R})$ mathematisch
durch ${\rm Erw}({\cal R})=\langle \psi|R|\psi\rangle$ gegeben ist. 

\subsection{Reine und gemischte Zust\"ande in der klassischen Mechanik} 

In der klassischen Mechanik der Punktteilchen ist ein reiner Zustand mathematisch durch einen
Punkt im Phasenraum -- dem Raum der m\"oglichen Orte und Impulse eines oder
mehrerer Teilchen -- gegeben. Wenn wir von einem Teilchen den Ort und seinen Impuls zu einem
Zeitpunkt kennen, haben wir das maximale Wissen, das wir \"uber dieses Teilchen haben k\"onnen.
Ort und Impuls legen als Anfangsbedingung eine Trajektorie im Phasenraum fest (hierf\"ur
verwendet man meist die Hamilton'sche Formulierung der Bewegungsgleichungen). Die Menge
der Observablen wird mathematisch durch die Menge der reell-wertigen Funktionen auf dem
Phasenraum dargestellt. Sei $f(x,p)$ eine Funktion auf dem Phasenraum, dann ist der
Funktionswert an einer bestimmten Stelle $(x_0,p_0)$ der Erwartungswert, den wir der Funktion
$f$ im Zustand $(x_0,p_0)$ zuordnen:
\begin{equation}
                         {\rm Erw:~} f \mapsto \langle f \rangle_{x_0,p_0} = f(x_0,p_0)  \, .
\end{equation}

Ein gemischter Zustand entspricht in der klassischen Physik meist einer teilweisen
Unkenntnis \"uber den Zustand eines Systems und wird mathematisch durch eine
Wahrscheinlichkeitsdichte $w(x.p)$ auf dem Phasenraum beschrieben. In diesem Fall ist
der Erwartungswert f\"ur eine Observable $f$ durch
\begin{equation}
             {\rm Erw:~} f \mapsto \langle f \rangle_{w} = \int_P f(x,p) w(x,p)\,{\rm d}x\,{\rm d}p   
\end{equation}
gegeben. Ein gemischter Zustand ordnet also jedem Punkt im Phasenraum (also jedem reinen
Zustand) eine gewisse Wahrscheinlichkeit zu und bildet dann den Mittelwert \"uber die Funktionswerte
an diesen Punkten. 

Ein Beispiel f\"ur einen klassischen Zustand ist die kanonische
Verteilung in der statistischen Mechanik: Sei $(x,p)$ der Phasenraum f\"ur ein System
von $N$ Teilchen (d.h., $x$ und $p$ sind jeweils $3N$-Vektoren) und sei $E(x,p)$ die Energie
dieser Teilchen im Zustand $(x,p)$, dann ist die Wahrscheinlichkeitsdichte des kanonischen
Zustands durch die Boltzmann-Verteilung
\begin{equation}
          w(x,p) = \frac{1}{Z} {\rm e}^{-\beta E(x,p)}
\end{equation}
gegeben. Hierbei ist $\beta=1/k_{\rm B}T$ die inverse Temperatur (mit der Boltzmann-Konstanten
$k_{\rm B}$) und $Z$ die Zustandssumme
\begin{equation}
          Z = \int_P {\rm e}^{-\beta E(x,p)}\,{\rm d}x\,{\rm d}p  \, ,
\end{equation}
die in diesem Fall die Wahrscheinlichkeitsdichte normiert. (Auf die Problematik, dass wegen der
Ununterscheidbarkeit gleichartiger Teilchen ein Faktor $N!$ sowie wegen des Ma\ss es 
im Phasenraum ein Faktor $h^N$ ber\"ucksichtigt werden m\"ussen, gehe ich hier nicht ein.) 

Der reine Zustand ist ein Spezialfall, bei dem die Wahrscheinlichkeitsdichte auf den Punkt
$(x_0,p_0)$ konzentriert ist: 
\begin{equation}
           w(x,p)_{x_0,p_0} = \delta(x-x_0) \delta(p-p_0) \, .
\end{equation}
In diesem Fall gilt f\"ur eine Observable $f$:
\begin{equation}
             \langle f \rangle = \int_P f(x,p) w(x,p)_{x_0,p_0}\,{\rm d}x\,{\rm d}p = f(x_0,p_0) \, .
\end{equation}

\subsection{Reine und gemischte Zust\"ande in der Quantenmechanik}

In der Quantenmechanik werden reine Zust\"ande mathematisch durch Strahlen
(1-dimensionale Untervektorr\"aume) in einem Hilbert-Raum (einem Vektorraum mit
Skalarprodukt) dargestellt. Meist w\"ahlen wir f\"ur diese Strahlen einen normierten
Vektor als Repr\"asentanten. In der Quantenmechanik von Punktteilchen ist dieser
normierte Vektor oft eine Wellenfunktion $\psi(x)$, die einen bestimmten Zustand (z.B.\ einen
bestimmten Energiezustand oder einen bestimmten Zustand zu anderen Quantenzahlen)
eines Teilchens beschreibt. Bei Systemen, die nur endlich viele Zust\"ande einnehmen
k\"onnen -- dazu z\"ahlen die Polarisationsfreiheitsgrade von Photonen oder die
Spinfreiheitsgrade von Elektonen oder Atomen -- schreiben wir diesen Vektor meist
als Ket-Vektor $|h\rangle$ (f\"ur ein Photon mit horizontaler Polarisierung), 
$|\uparrow_z\rangle$ (f\"ur ein Elektron mit Spin \glqq up\grqq\ bez\"uglich der $z$-Richtung) etc. 
Auch eine Wellenfunktion kann als Ket-Vektor geschrieben werden, allerdings geben
wir die Wellenfunktion meist durch ihr Argument an einem beliebigen Ort $x$ an, sodass
wir die Beziehung $\psi(x)= \langle x | \psi\rangle$ erhalten. 

Observable werden in der Quantentheorie mathematisch durch selbstadjungierte
lineare Operatoren auf dem Hilbert-Raum dargestellt. Handelt es sich bei dem Hilbert-Raum
um den Raum der quadratintegrierbaren Funktionen, also dem Raum der Wellenfunktionen,
so handelt es sich bei den selbstadjungierten Operatoren meist um Differentialoperatoren. 
Sei $A$ ein selbstadjungierter Operator auf einem Hilbert-Raum und $\psi$ ein reiner
Zustand, beschrieben durch einen normierten Ket-Vektor $|\psi\rangle$, dann ist das
Erwartungswertfunktional gegeben durch:
\begin{equation}
              A \rightarrow \langle A \rangle_\psi = \langle \psi |A| \psi \rangle  \, .
\end{equation}
F\"ur Wellenfunktionen $\psi(x)$ und lineare Differentialoperatoren $A(x,-{\rm i}\hbar\nabla)$
bedeutet dies:
\begin{equation}
       A \rightarrow \langle A \rangle_\psi 
                        = \int_{\mathbb{R}^3} \psi^*(x) A(x,-{\rm i}\hbar \nabla) \psi(x) \,{\rm d}x  \, .
\end{equation}

Ein reiner Zustand $\psi$ l\"asst sich auch durch einen Projektionsoperator $P_\psi$ auf den
1-dimensionalen Vektorraum, der diesem Zustand entspricht, darstellen. In diesem Fall
ist das Erwartungswertfunktional gegeben durch:
\begin{equation}
              A \rightarrow \langle A \rangle_\psi = {\rm Spur}\, P_\psi A  \, .
\end{equation}

Gemischte Zust\"ande werden in der Quantenmechanik ganz analog zur klassischen 
Mechanik formuliert: Jeder reine Zustand, beschrieben durch einen Projektionsoperator $P_i$, 
der mit einer Wahrscheinlichkeit $p_i$ in einem Ensemble vertreten ist, wird mit dieser
Wahrscheinlichkeit gewichtet und es wird \"uber alle reinen Zust\"ande summiert:
\begin{equation}
               \rho = \sum_i p_i P_i  \hspace{1cm}~~p_i\geq 0 ~,~ \sum_i p_i=1 \, .
\end{equation}
$\rho$ bezeichnet man als Dichtematrix. Ist eine solche Dichtematrix gegeben, ist das
Erwartungswertfunktional f\"ur einen Operator $A$ gegeben durch:
\begin{equation}
              A \rightarrow \langle A \rangle_\rho = {\rm Spur}\, \rho A   = \sum_i p_i \, {\rm Spur}\,P_iA \, .
\end{equation}
Auch f\"ur diesen Fall betrachten wir den kanonischen Zustand in der statistischen
Mechanik. Die Dichtematrix ist gegeben durch ($H$ ist der Hamilton-Operator, also der
Energieoperator in der Quantentheorie, und $E_i$ sind die Eigenwerte von $H$):
\begin{equation}
      \rho = \frac{1}{Z} {\rm e}^{-\beta H}  \hspace{1cm}  Z= {\rm Spur}\,{\rm e}^{-\beta H} =
                  \sum_i {\rm e}^{-\beta E_i} \, . 
\end{equation} 
Der Erwartungswert einer Observablen ist gegeben durch
\begin{equation}
               \langle A \rangle_\rho = \frac{1}{Z} {\rm Spur}\, {\rm e}^{-\beta H} A   
                        = \frac{1}{Z} \sum_i {\rm e}^{-\beta E_i} \, \langle E_i |A| E_i \rangle  \, .
\end{equation}

\section{Verschiedene Interpretationen des Zustandsbegriffs}

Bisher haben wir einen Zustand als \glqq Kenntnis \"uber ein System zur Vorhersage von 
Erwartungswerten\grqq\ bezeichnet. Physikalische realisiert ist ein Zustand durch ein Ensemble
gleichartig pr\"aparierter Systeme, mathematisch wird ein Zustand in der Quantentheorie durch
eine Dichtematrix -- ein reiner Zustand durch einen eindimensionalen Projektionsoperator bzw.\
einen Vektor (Wellenfunktion) als Repr\"asentanten -- dargestellt. Dies ist eine sehr \glqq vorsichtige\grqq\
Interpretation des Zustandsbegriffs, die in etwa der Kopenhagener Deutung entspricht. 

\subsection{Q-Bism -- Quanten-Bayesianismus}

Noch etwas extremer als die Kopenhagener-Deutung
in Richtung einer subjektiven Interpretation ist der Quanten-Bayesianismus oder kurz QBism 
(\cite{Fuchs} sowie die dortigen Literaturangaben), 
der unter einem Zustand ganz konkret \glqq meine Kenntnis\grqq\ \"uber ein System versteht. 
Damit ist man zwar auf der sicheren Seite, und -- wenn man ehrlich ist -- auf der Seite der
Praxis, denn mehr als meine Kenntnis kann ich zur Formulierung eines Zustands nicht
einbringen. Allerdings gibt man auch fast jede Form von Realismus auf. Das wird von den
Anh\"angern zwar bestritten, sie sprechen gerne von einem \glqq participatory realism\grqq, 
so wirklich \"uberzeugend ist dieser Realismus nicht. 

QBism f\"uhrt insbesondere zu folgender seltsamer Situation: Alice und Bob befinden sich
an zwei weit voneinander entfernten Orten und haben jeder in ihrem Labor
einen Teil eines verschr\"ankten Systems, beispielsweise ein Photon von einem im EPR-Zustand
verschr\"ankten Photonenpaar. Wenn Alice nun an ihrem Photon eine Messung vornimmt, erh\"alt
sie ein konkretes Messergebnis und reduziert entsprechend
den Zustand des Systems in einen separablen Zustand. Sie kann nun
vorhersagen, was Bob bei einer Messung an diesem Zustand entlang derselben Orientierung
als Ergebnis erhalten wird. Das bedeutet, Alice beschreibt das System nun durch einen
anderen Zustand als Bob, der nach wie vor den EPR-Zustand verwendet, solange er keine
weiteren Informationen hat. 

Damit erhebt sich die Frage, welchen Zustand die beiden Photonen \glqq wirklich\grqq\ haben.
F\"ur einen QBisten ist diese Frage sinnlos (wie \"ubrigens auch f\"ur andere, sogenannte
\glqq relative\grqq\ oder auch \glqq relationale\grqq\ Interpretationen). Ein Zustand ist
immer definiert in Bezug auf einen Beobachter; es handelt sich nicht um ein absolutes 
Objekt oder etwas, das
dem Quantensystem unabh\"angig vom Beobachter zugeschrieben werden kann.

\subsection{\glqq Neo-\grqq Kopenhagener Deutung}

W\"ahrend einige \"Au\ss erungen der Entwickler der Quantentheorie (beispielsweise
Bohr und Sommerfeld) darauf hindeuten, dass auch die urspr\"ungliche Kopenhagener
Deutung einen sehr subjektiven Zustandsbegriff hatte, denkt man heute bei der
Kopenhagener Deutung eher an eine objektivere Interpretation: Ein System wird
durch einen Zustand beschrieben, der sich auf das prinzipiell vorhandene Wissen
\"uber das System bezieht. Ob tats\"achlich ein Beobachter dieses Wissen hat, ist
dabei eher nebens\"achlich. Beispielsweise spielt es bei einem Interferenzexperiment
keine Rolle, ob wir tats\"achlich f\"ur jedes Quantenobjekt die \glqq Welcher Weg\grqq-Information
haben oder nicht: Wenn diese Information im Prinzip zug\"anglich ist, beobachtet man kein
Interferenzmuster. Es ist also keine subjektive Unkenntnis, die das Interferenzmuster
hervorruft, sondern eine prinzipielle Unm\"oglichkeit, die Kenntnis zu erlangen.

W\"ahrend sich im QBism der Zustand ausschlie\ss lich im Bewusstsein des Beobachters
befindet und keinerlei objektive Bedeutung hat, versucht man in der Neo-Kopenhagener Deutung
den Zustandsbegriff zu objektivieren und eher dem betrachteten System zuzuschreiben. 
Die Reduktion des Quantenzustands beruht im QBism ausschlie\ss lich auf dem
pers\"onlichen ge\"anderten Wissen \"uber ein System, w\"ahrend sie in der Kopenhagener
Deutung eher damit zusammenh\"angt, ob eine bestimmte Form von Information
prinzipiell vorhanden ist. 

\subsection{\glqq Viele Welten\grqq-Interpretation}

In der \glqq Viele Welten\grqq-Interpretation (bzw.\ \glqq Many-Worlds\grqq-Interpretation)
definiert der Quantenzustand $\psi$ die Realit\"at. Diese Realit\"at umfasst s\"amtliche
Universen, die seit dem Big Bang theoretisch m\"oglich sind. $\psi$ l\"asst sich beschreiben
als ein Funktional s\"amtlicher Teilchenorte (oder -- in einer Quantenfeldtheorie -- s\"amtlicher
Feldkonfigurationen) sowie s\"amtlicher Raumgeometrien, die sich aus dem Urknall h\"atten
entwickeln k\"onnen. Die zeitliche Entwicklung dieses Zustands folgt im Wesentlichen
einer Schr\"odinger-Gleichung (bzw.\ in einem allgemein relativistischen Rahmen einer
sogenannten Wheeler-de Witt-Gleichung). Es findet kein Kollaps (bzw.\ keine Reduktion)
der Wellenfunktion statt. Der Quantenzustand ist eine Superposition aller m\"oglichen
Anteile, die ein konsistentes Universum beschreiben, und in jedem Augenblick wird aus
jedem dieser Anteile wiederum eine Superposition von allen M\"oglichkeiten. 

Die Dekoh\"arenz zwischen den verschiedenen Anteilen in dieser Superposition
bez\"uglich der uns zur Verf\"ugung stehenden Observablen bewirkt, dass wir keine
Superpositionen oder Interferenzen zwischen den verschiedenen Anteilen wahrnehmen. Es ist
also unser subjektiver Eindruck, in einer ausgezeichneten Welt zu leben. 

Die Viele-Welten-Interpretation ist sehr verbreitet unter Physikerinnen und Physikern, die
sich mit Quantenkosmologie besch\"aftigen. Sie geht urspr\"unglich auf den Physiker
Hugh Everett zur\"uck, der eine Interpretation der Quantentheorie formulieren wollte,
bei der der Beobachter Teil des Quantensystems sein kann und dass keinen Bezug zu
einer \"ausseren klassischen Welt ben\"otigt \cite{Everett}. Die Bezeichnung Viele-Welten-Interpretation
wurde erst sp\"ater von Bryce deWitt gepr\"agt \cite{deWitt}. 

Die Viele-Welten-Interpretation ist eine rein deterministische Interpretation der Quantentheorie,
da es keinen Kollaps bzw.\ keine Reduktion der Wellenfunktion gibt. Au\ss erdem definiert diese
Interpretation eine -- wenn auch etwas seltsame -- Ontologie: Die objektive Realit\"at ist die
Wellenfunktion des Universums, die alle m\"oglichen Universen umfasst. 

\subsection{Bohm'sche Mechanik}

In der Bohm'schen Mechanik gibt es sowohl die Wellenfunktion $\psi(t)$ -- und hier ist
es von Bedeutung, dass wir von der Wellenfunktion im Ortsraum sprechen, d.h.\ der
Abbildung $\psi(t):\mathbb{R}^3 \rightarrow \mathbb{C}$ mit $x \mapsto \psi(x,t)$ -- als auch
das zugeh\"orige Teilchen, beschrieben durch seinen Ort $x(t)$. 
Das bedeutet, der momentane Zustand eines Systems ist 
durch das Paar $(\psi(t),x(t))$ gegeben. Der Ort $x(t)$ hat in der herk\"ommlichen 
Interpretation der Quantenmechanik
keine Bedeutung und entspricht einem sogenannten \glqq verborgenen Parameter\grqq. 

Das Feld $\psi(x,t)$ ist eine L\"osung der Schr\"odinger-Gleichung. Ist zu einem
Anfangszeitpunkt $t_0$ die statistische Verteilung der m\"oglichen Orte des Teilchens 
(d.h., die Wahrscheinlichkeit, das Teilchen bei einer Messung an einem Ort $x$ zu finden)
durch $|\psi(x,t_0)|^2$ gegeben, so bleibt dies auch f\"ur alle Zeiten so. Simulationen
solcher Systeme deuten darauf hin, dass eine anf\"angliche von $|\psi(x,t_0)|^2$
abweichende Verteilung im Allgemeinen sehr schnell gegen diese Verteilung konvergiert. 

In der Bohm'schen Mechanik gibt es keinen Kollaps der Wellenfunktion, aber
da sich das Teilchen an einem ganz bestimmten Ort befindet, kann man Teile
der Wellenfunktion, die f\"ur das weitere Verhalten des Teilchens nicht mehr relevant
sind, \glqq vergessen\grqq. Der genaue ontische Status des Feldes bleibt offen und hier
gibt es unterschiedliche Auffassungen.
F\"ur Bohm handelte es sich um ein F\"uhrungsfeld, das die Information dar\"uber
enth\"alt, wie sich das Teilchen bewegen soll. Manche sehen darin eine Art Potenzial.
In jedem Fall handelt es sich bei der Wellenfunktion um ein objektives Feld, dessen
Existenz unabh\"angig von einem Beobachter ist. Auch die Bohm'sche Mechanik
definiert eine Ontologie, d.h., sie macht keinen Bezug zu einer \glqq klassischen
Umgebung\grqq\ und der Messprozess kann unabh\"angig von einem Beobachter
innerhalb der Theorie interpretiert werden.

Einen objektiven Charakter hat allerdings nur die Wellenfunktion des gesamten
Kosmos. Wellenfunktionen f\"ur einzelne Teilchen, wie man sie oft in Publikationen
dargestellt sieht (z.B.\ f\"ur das Doppelspaltexperiment) sind Idealisierungen, sofern
die Teilchen nicht mit anderen Objekten verschr\"ankt sind. 

\begin{thebibliography}{99}
\bibitem{deWitt} deWitt, B.S., {\em Quantum Mechanics and Reality:
         Could the solution to the dilemma of indeterminism be a 
       universe in which all possible outcomes of an experiment
        actually occur?} Physics Today 23 (1970) 30--40.   
\bibitem{Everett} Everett, H., {\em Relative State Formulation of
         Quantum Mechanics}, Rev.\ Mod.\ Phys.\ {\bf 29} (1957) 454--462.\\
       {\em Theory of the Universal Wavefunction}; Thesis, Princeton
       University, Princeton (1956).       
\bibitem{Fuchs} Fuchs, Christopher A., Mermin, N.\ David, Schack, Ruediger; \textit{An Introduction
            to QBism with an Application to the Locality of Quantum Mechanics}; Am.\ J.\ Phys., Vol.\ 82,
            No.\ 8, 2014, 749--754.
\bibitem{Neumann} von Neumann, John; \textit{Mathematische Grundlagen der
     Quantenmechanik}, Verlag von Julius Springer, Berlin 1935. 
\bibitem{Schroedinger} Schr\"odinger, E.; {\em Die gegenw\"artige
        Situation in der Quantenmechanik}; Die Naturwissenschaften 23
        (1935) 807--812, 823--828,  844-849.      
\end{thebibliography}

\end{document}

